\section{Related work}
\label{sec:related}

\paragraph{Continuation and curricula.}
Classical continuation methods trace solutions through parameterized problems \citep{allgower1990continuation}. For nonconvex optimization, \citet{mobahi2015gaussian} analyze continuation through Gaussian-smoothed objectives, while \citet{hazan2016graduated} identify a class for which graduated optimization converges globally. Spatial filtering inside a network is not equivalent to smoothing its loss in parameter space, so these guarantees do not directly apply here. Related approaches to neural-network training include noisy activation and network deformations in Mollifying Networks \citep{gulcehre2016mollifying}, and curricula that progressively change which examples contribute to training \citep{bengio2009curriculum}. Our interventions keep the training examples and labels fixed and change their spatial processing. Their construction is motivated by continuation, without an established guarantee of easier optimization or minimizer tracking.
\paragraph{Spatial information during training.}
Curriculum by Smoothing \citep{sinha2020curriculum} is our direct starting point: Gaussian kernels filter feature maps, and their scale decreases during training. Changing image resolution is also an established training strategy. \citet{hoffer2019mixmatch} mix input sizes during training, while \citet{tan2021efficientnetv2} combine progressive image resizing with adaptive regularization in EfficientNetV2.

A particularly close precedent is SDPoint \citep{kuen2018stochastic}, which applies average-pooling downsampling at a randomly selected location and ratio during training. Its instances share parameters and provide different inference costs, with reported regularization benefits. Our resolution procedure instead fixes the insertion point within each run and restores spatial resolution according to a prescribed schedule. We examine this procedure alongside Gaussian filtering and their combination, with every selected method ending at the original inference architecture. This places our study within existing work on spatially varying computation, while focusing on the location and schedule of a temporary intervention.
\paragraph{Frequency and aliasing.}
\citet{rahaman2019spectral} study a bias towards low-frequency functions during neural-network learning. This is related to our motivation, but frequencies of the learned function in input space differ from spatial frequencies within a feature map. Downsampling also differs from low-pass filtering. \citet{zhang2019shift} introduce filtering before subsampling to improve shift consistency, and \citet{vasconcelos2021aliasing} study its placement and effects on generalization. Nonlinearities can introduce additional frequency components and aliasing \citep{michaeli2023aliasfree}. These works motivate attention to filter placement: a blur applied after a strided convolution does not undo aliasing already introduced by that convolution.
\paragraph{Implicit regularization and loss geometry.}
The selected procedures end in the same architecture, so a different final parameter count cannot explain their effects. Optimization and parameterization can nevertheless influence which functions are learned \citep{neyshabur2017geometry}. Early-stopping analyses characterize implicit regularization in least-squares models \citep{ali2019early}, while algorithmic stability provides another connection between optimization and generalization \citep{hardt2016stability}. These are possible frameworks for interpreting training--test differences, rather than established explanations of our observations. For empirical geometry, \citet{li2018visualizing} use filter-normalized directions and trajectory projections to visualize neural-network losses. Such views require care: reparameterization can change apparent sharpness without changing the predictor \citep{dinh2017sharp}, so a visually flatter surface alone does not establish better generalization.